\subsection{李群上的 Haar 测度}

设 G 是 n 维李群。于是 \(\bigwedge^{n}(\mathrm{T}_{e}(\mathrm{G}))\) 是 1 维的。因此（第13号），G 上 n 阶左不变微分形式的向量空间 S 是 1 维的。设 \((\omega_{1},\ldots,\omega_{n})\) 是 G 上 1 阶左不变微分形式空间的一组基；则 \(\omega_{1}\wedge\omega_{2}\wedge\cdots\wedge\omega_{n}\) 是 S 的一组基。

命题 54. 设 G 是 n 维李群，\(\omega\) 是 G 上的 n 阶左不变微分形式，\(\phi\) 是 G 的自同态。则

\[
\phi^ {*} (\omega) = (\det L (\phi)) \omega .
\]

记 \(\mathrm{L}(\phi) = u, \omega_e = f\) 且 \(\phi^*(\omega)_e = g\)。对任意 \(x_1, \ldots, x_n\)，它们属于 \(\mathrm{L}(\mathrm{G})\)，

\[
g (x _ {1}, \dots , x _ {n}) = f (u x _ {1}, \dots , u x _ {n}) = (\det u) f (x _ {1}, \dots , x _ {n})
\]

因此 \(\phi^{*}(\omega)_{e} = \det L(\phi). \omega_{e}\)。另一方面，若 \(g \in G\)，

\[
\phi \circ \gamma (g) = \gamma (\phi (g)) \circ \phi
\]

因此 \(\gamma(g)^{*}\phi^{*}(\omega)=\phi^{*}(\omega)\)。所以 \(\phi^{*}(\omega)\) 左不变，命题得证。

推论。对任意 \(g \in G\)，

\[
\delta (g) ^ {*} \omega = (\det \operatorname{Ad} g) \omega .
\]

\[
\delta (g) ^ {*} \omega = \delta (g) ^ {*} \gamma (g) ^ {*} \omega = (\operatorname{Int} g) ^ {*} \omega \text {   and   } L (\operatorname{Int} g) = \operatorname{Ad} g.
\]

设 G 是局部紧群，\(\phi\) 是 G 的自同态。假设存在 e 的开邻域 V、\(V'\)，使得 \(\phi(V) = V'\)，且 \(\phi \mid V\) 是从 G 到 G 的局部同构。设 \(\mu\) 是 G 上的左 Haar 测度。根据《积分学》第 VII 章 §1 命题9的推论，存在唯一的数 a \textgreater{} 0，使得 \(\phi(\mu \mid V) = a^{-1}\mu \mid V'\)。显然，a 与 V、\(V'\) 及 \(\mu\) 的选择无关。称其为 \(\phi\) 的模，记为 mod \(_{G}\) \(\phi\)，或简记为 mod \(\phi\)。当 \(\phi\) 是 G 的自同构时，便得到《积分学》第 VII 章 §1 的定义4。

命题 55. 假设 K 是局部紧的。设 \(\mu\) 是 K 的加法群上的 Haar 测度。设 G 是 \(n\) 维李群。

\begin{enumerate}
\def\labelenumi{(\roman{enumi})}
\item
  设 \(\omega\) 是 G 上非零的 \(n\) 阶左不变微分形式。则测度 mod \((\omega)_{\mu}\)（《可微与解析流形》，R，10.1.6）是 G 上的左 Haar 测度。若 K = R 且 G 具有由 \(\omega\) 定义的定向，则由 \(\omega\) 定义的测度（《可微与解析流形》，R，10.4.3）是 G 上的左 Haar 测度。
\item
  设 \(\phi\) 是 G 的 étale 自同态。则 mod \(\phi = \text{mod}\det L(\phi)\)。
\item
  显然。设 \(\mathrm{V}, \mathrm{V}'\) 是 \(e\) 的开邻域，使得 \(\phi(\mathrm{V}) = \mathrm{V}'\)，且 \(\phi \mid \mathrm{V}\) 是从 \(\mathrm{G}\) 到 \(\mathrm{G}\) 的局部同构。则
\end{enumerate}

\[
\begin{array}{r l r l} \phi^ {- 1} (\mathrm{mod} (\omega) _ {\mu}   |   \mathrm{V} ^ {\prime}) & = \mathrm{mod} (\phi^ {*} (\omega)) _ {\mu}   |   \mathrm{V}) & & \text {by transport of structure} \\ & = \mathrm{mod} (\det \mathrm{L} (\phi) \omega   |   \mathrm{V}) _ {\mu} & & \text {(Proposition 54)} \\ & = \mathrm{mod} \det \mathrm{L} (\phi) (\mathrm{mod} (\omega) _ {\mu}   |   \mathrm{V}) \end{array}
\]

因此 \(\phi = \mathrm{mod}\det \mathrm{L}(\phi)\)；这是由 \(\phi\) 的模的定义得到的。

推论。对任意 \(g \in \mathbf{G}\)，\(\Delta_{\mathbf{G}}(g) = (\text{mod det } \operatorname{Ad} g)^{-1}\)。特别地，\(\mathbf{G}\) 为幺模群的充要条件是 mod det \(\operatorname{Ad} g = 1\) 对所有 \(g \in \mathbf{G}\) 都成立。

\[
\begin{array}{r l} \Delta_ {\mathbf {G}} (g) & = (\text {mod} \operatorname{Int} g) ^ {- 1} \quad \text {(Integration, Chapter VII, § 1, formula (33))} \\ & = (\text {mod} \det \operatorname{L} (\operatorname{Int} g)) ^ {- 1} \quad \text {(Proposition 55)} \\ & = (\text {mod} \det \operatorname{Ad} g) ^ {- 1}. \end{array}
\]

注记。在保持命题52的假设和记号的情况下，假设 K 是局部紧的。设 \(\mu\) 是测度

\[
\operatorname{mod} \det (\mathrm{D} _ {n + i} m _ {k} (x _ {1}, \dots , x _ {n}, x _ {1}, \dots , x _ {n})) _ {1 \leqslant i, k \leqslant n} d x _ {1} \dots d x _ {n}
\]

在 \(\psi (\mathbf{U})\) 上。则 \(\psi^{-1}(\mu)\) 是 G 上某个 Haar 测度限制到 U 的结果。

命题 56. 设 G 是 n 维李群，H 是 p 维李子群，X 是李齐性空间 G/H。假设

\[
\det \operatorname{Ad} _ {\mathrm{L} (\mathrm{G})} h = \det \operatorname{Ad} _ {\mathrm{L} (\mathrm{H})} h
\]

对所有 \(h\in \mathbf{H}\) 均成立。则：

\begin{enumerate}
\def\labelenumi{(\roman{enumi})}
\item
  \(n - p\) 阶的 \(\mathbf{X}\) 上在 \(\mathbf{G}\) 下不变的微分形式都是解析的。
\item
  这些形式构成的向量空间是 1 维的。
\item
  若 \(\omega\) 是这样的非零形式且 K 是局部紧的，则 mod \((\omega)_{\mu}\) 是 X 上的非零测度，并且在 G 下不变。
\end{enumerate}

根据§ 1 第8号的例，\(\mathrm{Alt}^{n - p}(\mathrm{TX},\mathrm{K})\) 是解析向量 G-丛。设 \(x_0\) 是 \(e\) 在 \(\mathbf{X}\) 中的典范像；其稳定子为 H。\(\mathrm{Alt}^{n - p}(\mathrm{TX},\mathrm{K})\) 在 \(x_0\) 处的纤维是 \(\bigwedge^{n - p}\mathrm{T}_{x_0}(\mathbf{X})^*\)，而 \(\mathrm{T}_{x_0}(\mathbf{X})\) 与 \(\mathrm{L}(G) / L(\mathrm{H})\) 典范地识别。若 \(h\in \mathrm{H}\)，\(\tau_h\) 是 \(\mathbf{X}\) 的自同构，由 \(h\) 定义；在取商时，它由 G 的自同构 \(g\mapsto hgh^{-1}\) 导出。因此，自同构 \(\mathrm{T}_{x_0}(\tau_h)\) 在取商时由 \(\mathrm{Ad}_{\mathrm{L(G)}}(h)\) 导出。由于

\[
\det \mathrm{Ad} _ {\mathrm{L(G)}} h = (\det \mathrm{Ad} _ {\mathrm{L(H)}} h). (\det \mathrm{T} _ {x _ {0}} (\tau_ {h})),
\]

假设推出 \(\det \mathrm{T}_{x_0}(\tau_h) = 1\)。所以 \(\bigwedge^{n - p}\mathrm{T}_{x_0}(\mathbf{X})^*\) 的每个元素都在 H 下不变。于是 (i) 和 (ii) 由§ 1 第8号命题17的推论1推出，而 (iii) 显然。

由《积分学》第 VII 章 § 2 定理3的推论2，存在一个在 G 下不变的 X 上非零正测度；因为命题56的假设推出 \(\Delta_{\mathbf{G}}|H = \Delta_{\mathbf{H}}\)（命题55的推论）。

命题 57. 设 G 是 n 维李群。~取 \(\bigwedge^{n}\mathrm{T}_{e}(\mathrm{G})^{*}\) 的一组基；借助 \(\bigwedge^{n}\mathrm{T}(\mathrm{G})^{*}\) 的右（分别为左）平凡化，我们可以将这个向量丛与平凡向量丛 \(\mathbf{G} \times \mathbf{K}\) 识别，从而标量微分算子的转置被识别为标量微分算子。

于是，若 \(u \in \mathrm{U}(\mathrm{G})\)，\(\mathrm{L}_u\)（分别为 \(\mathbf{R}_u\)）的转置是 \(\mathrm{L}_u^\vee\)（分别为 \(\mathbf{R}_u^\vee\)）。

我们考虑将 \(\bigwedge^{n}\mathrm{T}(\mathrm{G})^{*}\) 用右不变形式 \(\omega\) 平凡化的情形。

假设对于 U(G) 中的元素 \(u_{1}, u_{2}\) 命题已经得到证明。则

\[
\begin{array}{l l} ^ {t} (\mathrm{L} _ {u _ {1} * u _ {2}}) = ^ {t} (\mathrm{L} _ {u _ {1}} \circ \mathrm{L} _ {u _ {2}}) & \text {(Proposition 23)} \\ = ^ {t} (\mathrm{L} _ {u _ {2}}) \circ^ {t} (\mathrm{L} _ {u _ {1}}) & \text {(Diff. \& Anal. Man., R, 14.3.3)} \\ = \mathrm{L} _ {u _ {2}} ^ {\vee} \circ \mathrm{L} _ {u _ {1}} ^ {\vee} & \text {by hypothesis} \\ = \mathrm{L} _ {u _ {2} * u _ {1}} ^ {\vee \vee} & \text {(Proposition 23)} \\ = \mathrm{L} _ {(u _ {1} * u _ {2}) \vee} & \text {(Proposition 7)} \end{array}
\]

因此命题对 \(u_{1} * u_{2}\) 成立。所以只需在 \(u \in \mathrm{T}_{e}(\mathrm{G})\) 时证明命题。现在，\(\mathbf{L}_u\) 由 G 在 G 上的右作用定义（第6号），因而 \(\theta_{\mathrm{L}_u}\omega = 0\)，因为 \(\omega\) 右不变（《可微与解析流形》，R，8.4.5）；所以，若 \(f\) 是定义在 \(e\) 的一个开邻域上且取值于 K 的解析函数，则 \(\theta_{\mathrm{L}_u}(f\omega) = (\theta_{\mathrm{L}_u}f)\omega\)（《可微与解析流形》，R，8.4.8）。利用上述识别以及《可微与解析流形》R，14.4.1，\(\mathbf{L}_u\) 的转置是 \(-\mathbf{L}_u\)，即 \(\mathbf{L}_u^\vee\)。

推论。设 G 是有限维实李群，\(\mu\)（分别为 \(\nu\)）是 G 上的左（分别为右）Haar 测度，k 是满足 \(\geqslant0\) 的整数，\(u\in\mathrm{U}_{k}(\mathrm{G})\)，f 和 g 是 G 上具有紧支集的 \(C^{k}\) 类实值函数。则

\[
\int_ {G} \left(R _ {u} f\right) g d \mu = \int_ {G} f (R _ {u} ^ {\vee} g) d \mu
\]

\[
\int_ {\mathbf {G}} \left(\mathrm{L} _ {u} f\right) g d \nu = \int_ {\mathbf {G}} f (\mathrm{L} _ {u} ^ {\vee} g) d \nu .
\]

这由命题57及《可微与解析流形》R，14.3.8推出。

